Source : Eurostat, comptes annuels des administrations publiques (gov\_10a\_main, gov\_10a\_taxag, gov\_10a\_exp), PIB (nama\_10\_gdp) ; Insee, dette trimestrielle de Maastricht par sous-secteur (contributions consolidées, 4e trimestre). Calculs de l'auteur (dette/attribution.py).\par Lecture : la dette créée depuis 1995 (1 355 Md\euro{} fin 2025) est la somme des déficits primaires de 1996 à 2025 et des intérêts qu'ils ont générés ; le reste de la dette est le stock de 1995, ses intérêts et les ajustements stock-flux, laissés à l'émetteur. Une barre qui dépasse à gauche de zéro signale une contribution négative (excédents primaires cumulés). La méthode par fonction répartit la dette entre fonctions et non entre administrations : elle compte en santé et protection sociale le RSA versé par les départements ou les aides au logement.
